\documentclass[11pt,a4paper]{altacv}

\geometry{left=1.5cm,right=1.5cm,top=1.cm,bottom=1.2cm,columnsep=1.5cm}

\usepackage[utf8]{inputenc}
\usepackage[T1]{fontenc}
\usepackage[french]{babel}
\usepackage{paracol}
\usepackage{xcolor}
\usepackage{hyperref}
\usepackage{fontawesome5}
\usepackage{setspace}

\definecolor{SlateGrey}{HTML}{2E2E2E}
\definecolor{LightGrey}{HTML}{666666}
\definecolor{DarkPastelRed}{HTML}{8AC0D9}
\definecolor{PastelRed}{HTML}{00B0DA}
\definecolor{GoldenEarth}{HTML}{B4AD0C}
\colorlet{name}{black}
\colorlet{tagline}{PastelRed}
\colorlet{heading}{DarkPastelRed}
\colorlet{headingrule}{GoldenEarth}
\colorlet{subheading}{PastelRed}
\colorlet{accent}{PastelRed}
\colorlet{emphasis}{SlateGrey}
\colorlet{body}{LightGrey}

\renewcommand{\familydefault}{\sfdefault}

\name{\Large Guillaume BOILEAU}
\tagline{Chef de projet technique SI | SPOC, besoins techniques, build/run, risques, cycle en V \& coordination DSI}
\personalinfo{
  \email{guillaume.boileaupro@gmail.com}
  \phone{+33 6 59 94 85 84}
  \location{Alpes-Maritimes, France}
  \linkedin{guillaume-boileau-phd-891b81265}
}

\setlength{\parskip}{0.75\baselineskip}

\begin{document}
\makecvheader
\vspace{-0.6cm}

\cvsection{Lettre de motivation}
\vspace{-0.4cm}

{\color{accent}\textbf{Objet :}} \textit{Candidature -- Chef de projet stratégique informatique -- Valbonne}

{\color{body}

Madame, Monsieur,

Je vous adresse ma candidature pour le poste de \textbf{Chef de projet stratégique informatique}. Votre offre m'intéresse particulièrement car elle correspond à un rôle de coordination technique à forte valeur ajoutée : être point d'entrée des demandes, qualifier les besoins, consolider les contraintes, fédérer les parties prenantes, identifier les risques, préparer les arbitrages et assurer un reporting fiable entre les équipes de la DSI.

Mon parcours s'est construit à l'interface entre \textbf{développement logiciel, systèmes complexes, validation, support opérationnel, documentation et coordination transverse}. J'ai travaillé dans des environnements où il faut comprendre rapidement les enjeux techniques, rendre visibles les dépendances, clarifier les jalons, structurer les actions et maintenir un niveau d'information partagé entre plusieurs contributeurs.

Au \textbf{CNES}, j'ai contribué au segment sol de la mission spatiale \textbf{LISA}, dans un environnement CNES/ESA associant chaînes de données mission L0--L3, cycle en V, intégration de modèles, validation technique et revues collaboratives. J'y ai développé \textbf{CATpyLISA}, une plateforme Python permettant de structurer des workflows, comparer des modèles, qualifier des écarts, valider des résultats et produire une documentation exploitable par différents contributeurs.

Cette expérience est directement transposable aux missions de la Direction Technique de la CNAF : qualifier une demande technique, identifier les contraintes, consolider les dépendances, produire une fiche claire de besoin, suivre les actions, alerter sur les risques et préparer les éléments utiles aux comités. Elle m'a également donné une pratique solide du cycle en V, de l'IVVQ, de la traçabilité, de la documentation et du transfert de connaissances vers les équipes en charge du pilotage opérationnel.

Plus récemment, chez \textbf{VEV Platform Services France}, j'ai travaillé sur une plateforme logicielle industrielle connectée à des infrastructures de recharge électrique. Cette expérience m'a permis de renforcer ma compréhension des enjeux build/run : intégration logicielle, flux de données opérationnels, support technique, analyse d'incidents, continuité de service, diagnostic, formalisation de constats et contribution à la fiabilisation de services numériques en exploitation.

Pour la CNAF, je peux apporter un profil capable de jouer un rôle de \textbf{SPOC technique} entre directions et équipes techniques : écouter les besoins, consolider les informations, coordonner les acteurs, suivre les actions, identifier les risques, déclencher les alertes et produire une synthèse claire pour faciliter les décisions. Je suis habitué à gérer plusieurs sujets simultanément, à travailler avec rigueur, à documenter les échanges et à communiquer positivement dans des contextes complexes.

Je dispose également d'une culture SI et logicielle utile pour dialoguer avec des architectes, des équipes build, run, production, support et projet : Python, TypeScript, Angular, Node.js, Linux, Git, Docker, CI/CD, monitoring/logs, OpenSearch, Sentry, Confluence, Jira et Zenhub. Sans me présenter comme expert infrastructure DSI, je possède le socle technique nécessaire pour comprendre les enjeux, poser les bonnes questions et faciliter la coordination.

Le rôle proposé par la CNAF représente pour moi une continuité logique vers un poste de pilotage technique stratégique, au service d'une DSI à fort impact public. La mission de la CNAF, la dimension de service aux équipes de développement, la standardisation et l'automatisation de la plateforme technique correspondent à mon envie de contribuer à des projets utiles, structurants et exigeants.

Je serais heureux de pouvoir échanger avec vous afin de vous présenter plus en détail mon parcours et la manière dont mon expérience en coordination technique, validation, support run, documentation et analyse de risques peut contribuer efficacement aux missions de la Direction Technique de la CNAF.

Je vous prie d'agréer, Madame, Monsieur, l'expression de mes salutations distinguées.

\textbf{Guillaume Boileau}

}

\end{document}
